\chapter{Matériel programmable I~: architecture et chaîne d'outils}\label{ch:nw:programmable-hardware-i-architecture-and-toolchain}

Une liaison Ethernet à \qty{10}{\giga\bit\per\second} livre à un circuit ses octets au rythme de dix milliards de bits par seconde~; un circuit qui les prend
64 bits à la fois doit accepter un mot toutes les \qty{6.4}{\nano\second}, à 156.25~MHz. Un circuit construit de registres et de logique entre eux
met un nombre fixe de ces cycles pour décider, et il met ce même nombre à chaque fois~: aucun défaut de cache (cache miss), aucune interruption, aucun ordonnanceur, aucun
ramasse-miettes (garbage collector) ne s'intercale entre le dernier octet d'un cours et la décision. Le déterminisme, plutôt que la vitesse brute, est ce que le matériel programmable
vend à une société de trading, et il s'achète par une manière différente d'écrire des programmes.

Ce chapitre présente les dispositifs, les langages et les outils, et écrit un premier circuit~: un extracteur de champ qui prélève un cours
dans une trame pendant que la trame arrive encore, simulé cycle par cycle dans deux simulateurs indépendants et vérifié contre un
modèle Python. Le chapitre~7 construit les circuits de trading par-dessus celui-ci.

\section{Le tissu~: blocs logiques, mémoires, émetteurs-récepteurs et horloges}

\begin{definition}[Circuit logique programmable, bloc logique configurable]\label{def:nw:programmable-hardware-i-architecture-and-toolchain:clb}
Un \emph{circuit logique programmable}\index{circuit logique programmable} (FPGA) est une puce dont la logique et le câblage sont configurés
après la fabrication par le chargement d'un fichier de configuration~: un réseau de petits éléments logiques programmables, de mémoires, de blocs arithmétiques et
d'émetteurs-récepteurs série, reliés par un routage programmable. Un \emph{bloc logique configurable}\index{bloc logique configurable} est l'élément logique répété du réseau~:
un groupe de tables de correspondance, dont chacune calcule n'importe quelle fonction booléenne de quelques entrées en la lisant
dans une petite mémoire, avec des bascules pour enregistrer leurs sorties et une chaîne de retenue pour l'arithmétique.
\end{definition}

\begin{definition}[Block RAM]\label{def:nw:programmable-hardware-i-architecture-and-toolchain:bram}
Une \emph{block RAM}\index{block RAM} est une mémoire dédiée embarquée sur la puce, d'une capacité de quelques dizaines de kilobits, dont un FPGA compte des centaines ou
des milliers, chacune lisible et inscriptible en un cycle d'horloge, typiquement par deux ports indépendants.
\end{definition}

Les ressources sont concrètes. Dans la famille UltraScale d'AMD, chaque \omterm{def:nw:programmable-hardware-i-architecture-and-toolchain:clb}{bloc logique configurable} contient huit tables de correspondance à 6~entrées et seize
bascules, avec une chaîne de retenue de 8~bits, et chaque \omterm{def:nw:programmable-hardware-i-architecture-and-toolchain:bram}{block RAM} contient 36 kilobits, utilisables comme une mémoire ou comme deux moitiés indépendantes. Une table de correspondance à 6~entrées
calcule n'importe quelle fonction de six bits en un niveau de logique~; une fonction de davantage de bits demande plusieurs tables sur plusieurs niveaux,
et la profondeur de cet arbre de niveaux, avec le câblage entre les tables, fixe la fréquence à laquelle le circuit peut être cadencé. À côté de la logique
se trouvent les émetteurs-récepteurs série, qui transforment le flux de bits du module optique en mots parallèles, et les ressources d'horloge, qui
distribuent les horloges et permettent à des parties du circuit de fonctionner à des cadences différentes.

\begin{omfigure}
\begin{tikzpicture}[font=\footnotesize,b/.style={draw,rounded corners=2pt,minimum height=0.6cm,align=center,inner sep=3pt}]
\foreach \i in {0,...,5} {\foreach \j in {0,1,2} {\node[b,fill=omDef!12,minimum width=0.9cm] at (1.1*\i,0.9*\j) {CLB};}}
\foreach \j in {0,1,2} {\node[b,fill=omMeth!20,minimum width=0.9cm] at (6.9,0.9*\j) {BRAM};}
\foreach \j in {0,1,2} {\node[b,fill=omThm!15,minimum width=0.9cm] at (8.25,0.9*\j) {DSP};}
\node[b,fill=gray!15,minimum height=2.5cm,minimum width=1.1cm] at (-1.5,0.9) {émetteurs-\\ récepteurs\\ série};
\node[b,fill=gray!8,minimum width=9.6cm] at (3.25,-0.95) {horloge et routage entre tous les blocs};
\node[b,fill=white,align=left,anchor=west] at (9.2,0.9) {un CLB~: 8 tables\\ de correspondance à 6 entrées,\\ 16 bascules,\\ une chaîne de retenue};
\end{tikzpicture}

\omcaption{Le tissu d'un FPGA, schématiquement~: des colonnes de \omterm{def:nw:programmable-hardware-i-architecture-and-toolchain:clb}{blocs logiques configurables}, de \omterm{def:nw:programmable-hardware-i-architecture-and-toolchain:bram}{block RAM} et de blocs arithmétiques (DSP), des
émetteurs-récepteurs série en bordure qui transforment le flux optique en mots, et l'horloge et le routage qui les relient. Les circuits réels
comptent des centaines de milliers de tables de correspondance~; le contenu des CLB est celui de la famille UltraScale d'AMD.}\label{fig:nw:programmable-hardware-i-architecture-and-toolchain:fabric}
\end{omfigure}

\section{Décrire le matériel~: niveau transfert de registres et synthèse de haut niveau}

\begin{definition}[Niveau transfert de registres, langage de description de matériel]\label{def:nw:programmable-hardware-i-architecture-and-toolchain:rtl}
Un circuit écrit au \emph{niveau transfert de registres}\index{niveau transfert de registres} (RTL) décrit, pour chaque cycle d'horloge, comment les valeurs
contenues dans les registres sont calculées à partir des valeurs précédentes des registres et des entrées, par une logique combinatoire. Un \emph{langage de
description de matériel}\index{langage de description de matériel} (Verilog, SystemVerilog, VHDL) est un langage dans lequel de tels circuits sont écrits~:
ses instructions décrivent des circuits qui fonctionnent tous à la fois, et non des instructions exécutées l'une après l'autre.
\end{definition}

\begin{definition}[Synthèse logique, placement-routage]\label{def:nw:programmable-hardware-i-architecture-and-toolchain:synth}
La \emph{synthèse logique}\index{synthèse logique} traduit une description RTL en un réseau de primitives du dispositif (tables de correspondance,
bascules, mémoires). Le \emph{placement-routage}\index{placement-routage} (place and route) affecte chaque primitive à un emplacement de la puce et choisit les
fils programmables qui les relient~; son résultat est le fichier de configuration chargé dans le dispositif.
\end{definition}

\begin{definition}[Synthèse de haut niveau]\label{def:nw:programmable-hardware-i-architecture-and-toolchain:hls}
La \emph{synthèse de haut niveau}\index{synthèse de haut niveau} (HLS) compile une fonction écrite en C ou en C++ en un circuit RTL, en choisissant
ses registres et son ordonnancement~; le concepteur la guide par des directives, telles que l'intervalle d'initiation du pipeline, c'est-à-dire le nombre de cycles entre
deux entrées successives (un intervalle de un accepte une nouvelle entrée à chaque cycle).
\end{definition}

La différence de raisonnement est la première leçon du chapitre. Une boucle C++ sur les huit octets d'un mot s'exécute huit fois, l'une après
l'autre~; la même boucle en SystemVerilog décrit huit copies de son corps, qui existent toutes et calculent toutes au même cycle.
Le \Cref{lst:nw:programmable-hardware-i-architecture-and-toolchain:field} est le cœur combinatoire de l'extracteur de champ~: pour chacun des
huit voies d'octets du battement, il demande si cet octet appartient au champ et, dans ce cas, le décale dans la valeur, le tout en un cycle. Ce qui serait
un branchement conditionnel dépendant des données dans un logiciel devient un ensemble de multiplexeurs dont le délai ne dépend pas des données.

\omcode{code/firm/hdlkit/hdl/hdk_field.sv}{19}{29}{Logique combinatoire de l'extracteur de champ en SystemVerilog~: huit voies d'octets examinées
dans le même cycle~; la boucle décrit huit copies du matériel, et non huit itérations.}\label{lst:nw:programmable-hardware-i-architecture-and-toolchain:field}

\section{Fermeture temporelle~: périodes d'horloge, chemins critiques et pipelinage}

\begin{definition}[Chemin critique, fermeture temporelle]\label{def:nw:programmable-hardware-i-architecture-and-toolchain:timing}
Le \emph{chemin critique}\index{chemin critique} d'un circuit synchrone est son chemin le plus lent d'un registre au suivant~: le délai
de sa logique et de son câblage, auquel s'ajoute la surcharge propre des registres, ce qui fixe la période d'horloge la plus courte que le circuit supporte. La \emph{fermeture
temporelle}\index{fermeture temporelle} (timing closure) est le travail qui consiste à faire tenir chaque chemin entre registres dans la période d'horloge visée.
\end{definition}

\begin{definition}[Franchissement de domaine d'horloge]\label{def:nw:programmable-hardware-i-architecture-and-toolchain:cdc}
Un \emph{franchissement de domaine d'horloge}\index{franchissement de domaine d'horloge} est un signal qui passe entre des parties d'un circuit cadencées par des horloges sans lien entre elles~;
il doit passer par un synchroniseur ou une FIFO asynchrone, sinon il sera parfois échantillonné pendant qu'il change.
\end{definition}

\begin{proposition}[Ce que le pipelinage apporte]\label{prop:nw:programmable-hardware-i-architecture-and-toolchain:pipe}
Supposons qu'une décision demande $L$ niveaux de logique, de délai $t_{\mathrm{lvl}}$ chacun, et qu'un registre ajoute $t_{\mathrm{reg}}$. Découpé en $s$
étages enregistrés comportant chacun au plus $\lceil L/s\rceil$ niveaux, le circuit a pour période d'horloge $T_{\mathrm{clk}} = \lceil L/s\rceil\,t_{\mathrm{lvl}}
+ t_{\mathrm{reg}}$ et pour latence $s\,T_{\mathrm{clk}}$. Le pipelinage élève la fréquence d'horloge maximale $f_{\mathrm{clk}} = 1/T_{\mathrm{clk}}$
et ne réduit jamais la latence~: $s\,T_{\mathrm{clk}} \ge L\,t_{\mathrm{lvl}} + s\,t_{\mathrm{reg}}$.
\end{proposition}

\begin{proof}
Chaque étage est un chemin registre à registre comportant au plus $\lceil L/s\rceil$ niveaux, donc la période est son délai. La décision traverse $s$
étages, une période chacun, et $s\lceil L/s\rceil \ge L$.
\end{proof}

La proposition explique un choix de conception qui paraît à rebours~: un circuit de trading ajoute des registres, donc des cycles, pour aller plus vite. La liaison fixe
l'horloge. Une liaison à \qty{10}{\giga\bit\per\second} prise 64 bits à la fois demande 156.25~MHz~; une liaison à \qty{25}{\giga\bit\per\second} à la
même largeur demande 390.625~MHz. Une décision qui tient en un cycle à la première horloge peut en demander trois à la seconde, et sa latence en nanosecondes
augmente alors un peu, tandis que le circuit suit une liaison deux fois et demie plus rapide. La \Cref{fig:nw:programmable-hardware-i-architecture-and-toolchain:pipeline}
illustre le compromis avec des délais de modèle de \qty{0.5}{\nano\second} par niveau et \qty{0.4}{\nano\second} par registre, présentés comme des valeurs de modèle~:
les délais d'un dispositif réel proviennent de ses rapports de temporisation.

\begin{omfigure}
\begin{tikzpicture}
\begin{groupplot}[group style={group size=2 by 1,horizontal sep=1.8cm},width=6.6cm,height=5.0cm,xmin=0.5,xmax=6.5,
  xtick={1,2,3,4,5,6},xlabel={étages registrés},tick label style={font=\small},label style={font=\small},grid=major,
  grid style={gray!25},table/col sep=comma]
\nextgroupplot[ymin=0,ymax=800,ylabel={horloge maximale (MHz)}]
\addplot[omDef,very thick,mark=*,mark size=1.5pt] table[x=stages,y=fmax_mhz]{figdata/networks/06-programmable-hardware-i-architecture-and-toolchain/pipeline.csv};
\addplot[gray,dashed,domain=0.5:6.5] {156.25};
\addplot[omThm,dashed,domain=0.5:6.5] {390.625};
\node[font=\scriptsize,anchor=south west] at (axis cs:3.6,160) {10G à 64 bits};
\node[font=\scriptsize,anchor=north west,omThm] at (axis cs:4.1,385) {25G à 64 bits};
\nextgroupplot[ymin=0,ymax=10,ylabel={latence à sa propre horloge (ns)}]
\addplot[omMeth,very thick,mark=square*,mark size=1.5pt] table[x=stages,y=latency_ns]{figdata/networks/06-programmable-hardware-i-architecture-and-toolchain/pipeline.csv};
\end{groupplot}
\end{tikzpicture}

\omcaption{Modèle~: une décision de 12 niveaux de logique découpée en étages registrés (\qty{0.5}{\nano\second} par niveau,
\qty{0.4}{\nano\second} par registre, valeurs de modèle). À gauche~: l'horloge maximale, à comparer aux horloges qu'exigent les \omterm{def:nw:networking-for-trading:ser}{débits de ligne}~; à droite~: la latence
à cette horloge. Cinq étages ne font pas mieux que quatre, car $\lceil 12/5\rceil = \lceil 12/4\rceil = 3$. Données~:
\texttt{fig\_fpga.py}.}\label{fig:nw:programmable-hardware-i-architecture-and-toolchain:pipeline}
\end{omfigure}

\omcode{code/firm/hdlkit/hdl/hdk_cmp.sv}{13}{32}{Une comparaison en un étage et en deux~: la version à deux étages compare quatre tranches
de 16 bits, enregistre les résultats, puis les combine, un cycle plus tard.}\label{lst:nw:programmable-hardware-i-architecture-and-toolchain:cmp}

\section{Simulation et vérification}

\begin{definition}[Banc de test, modèle au cycle près]\label{def:nw:programmable-hardware-i-architecture-and-toolchain:tb}
Un \emph{banc de test}\index{banc de test} (testbench) est le code qui pilote les entrées d'une conception en simulation et vérifie ses sorties. Un \emph{modèle au cycle
près}\index{modèle au cycle près} (cycle-accurate model) est un modèle logiciel d'une conception qui reproduit ses sorties cycle par cycle, utilisé comme référence
à laquelle le matériel est vérifié, et comme substitut rapide de celui-ci.
\end{definition}

Une conception est vérifiée bien avant d'être synthétisée, car une conception synthétisée est lente à construire et difficile à observer. Ce livre
simule chaque conception avec deux outils indépendants. Verilator compile le SystemVerilog en un modèle C++ piloté par un \omterm{def:nw:programmable-hardware-i-architecture-and-toolchain:tb}{banc de test} en C++~: c'est
un compilateur plutôt qu'un simulateur traditionnel, et il est rapide. Icarus Verilog interprète la même source avec un \omterm{def:nw:programmable-hardware-i-architecture-and-toolchain:tb}{banc de test} SystemVerilog. Les
deux \omterm{def:nw:programmable-hardware-i-architecture-and-toolchain:tb}{bancs de test} lisent le même stimulus et impriment les mêmes lignes d'événements, et le test les compare cycle par cycle et à un modèle de référence en Python~;
un désaccord entre les simulateurs est un bogue de la conception (le plus souvent une ambiguïté que le langage autorise) avant d'être un bogue
de l'un des outils. Les conceptions sont écrites pour les versions installées ici (Verilator 4.038, Icarus Verilog 11) et pour leurs versions plus récentes,
et les tests échouent, au lieu d'être ignorés, lorsqu'un simulateur manque.

\begin{table}[ht]
\centering\small
\begin{tabular}{rll}
\toprule
Cycle & Entrée acceptée & Sorties enregistrées à la fin du cycle \\
\midrule
0 & mot 0~: octets \texttt{10}--\texttt{17} & --- \\
1 & mot 1~: octets \texttt{18}--\texttt{1f}, dernier & --- \\
2 & --- & champ \texttt{16171819}~; CRC-32 \texttt{f4a7fd67} \\
3 & --- & comparaison (un étage)~: au plus le seuil \\
4 & --- & comparaison (deux étages)~: au plus le seuil \\
\bottomrule
\end{tabular}
\caption{La première trame passant par \texttt{hdk\_top} dans Verilator, sans contre-pression~: une trame de 16 octets, allant des octets \texttt{10} à \texttt{1f}, le
champ aux octets 6 à 9, le seuil égal au champ. Le tampon de débordement (skid buffer) enregistre chaque beat accepté~; le champ et le CRC sont enregistrés
au cycle qui suit le beat qui les complète. Icarus affiche les mêmes cycles. Données~: \texttt{nw\_fpga.first\_frame}.}\label{tab:nw:programmable-hardware-i-architecture-and-toolchain:cycles}
\end{table}

\begin{omfigure}
\begin{tikzpicture}[font=\footnotesize,b/.style={draw,rounded corners=2pt,minimum height=0.7cm,align=center,inner sep=3pt},>=Stealth]
\node[b,fill=omDef!8] (hls) at (-3.4,0) {C++ avec\\ directives HLS};
\node[b,fill=omDef!15] (rtl) at (0,0) {RTL\\ (SystemVerilog)};
\node[b,fill=omMeth!18] (syn) at (2.8,0) {synthèse};
\node[b,fill=omMeth!18] (par) at (5.1,0) {placement-\\ routage};
\node[b,fill=omMeth!18] (sta) at (7.5,0) {analyse\\ temporelle};
\node[b,fill=gray!12] (bit) at (10.2,0) {fichier de\\ configuration};
\node[b,fill=omThm!12] (sim) at (4.6,-1.7) {simulation~: Verilator, Icarus, face à un modèle de référence};
\draw[->] (hls) -- node[above,font=\scriptsize]{HLS} (rtl);
\draw[->] (rtl) -- (syn);
\draw[->] (syn) -- (par);
\draw[->] (par) -- (sta);
\draw[->] (sta) -- node[above,font=\scriptsize]{tenue} (bit);
\draw[->,omThm] (sta.north) -- ++(0,0.5) -| node[pos=0.25,above,font=\scriptsize]{temporisation non atteinte~: restructurer, pipeliner} (rtl.north);
\draw[<->,omThm] (rtl.south) |- (sim.west);
\end{tikzpicture}

\omcaption{La chaîne d'outils. La conception est simulée face à son modèle de référence avant la synthèse~; la synthèse, le placement et le routage produisent une
conception dont la temporisation est analysée~; un chemin qui n'atteint pas l'horloge renvoie le concepteur au RTL.}\label{fig:nw:programmable-hardware-i-architecture-and-toolchain:flow}
\end{omfigure}

\section{Cartes réseau intelligentes}

\begin{definition}[SmartNIC]\label{def:nw:programmable-hardware-i-architecture-and-toolchain:smartnic}
Une \emph{SmartNIC}\index{SmartNIC} est une \omterm{def:nw:networking-for-trading:frame}{carte réseau} (NIC) dotée de son propre traitement programmable (un FPGA, ou des processeurs) entre les
ports réseau et l'hôte, de sorte que les paquets puissent être analysés, filtrés, horodatés, acquittés ou générés sur la carte sans atteindre
l'hôte.
\end{definition}

Pour une société de trading, la carte FPGA est le point où le matériel rejoint le reste du système~: ses émetteurs-récepteurs terminent les liaisons de la bourse, sa
logique décode et décide, et son bus la relie aux logiciels de l'hôte, qui la configurent et reçoivent ce qu'elle a vu. Une conception
sur une telle carte lit la trame à mesure qu'elle arrive, avant que le dernier octet ne soit reçu, ce que le tableau de ce chapitre montre à petite échelle~:
le champ est connu au cycle qui suit le beat qui le complète, alors que le reste de la trame est encore à venir. Le chapitre~7 décide à partir de lui.

\section{Tutoriel~: un extracteur de champ dans deux simulateurs}\label{tut:nw:programmable-hardware-i-architecture-and-toolchain:tut}

\begin{tutorial}
\textbf{Objectif.} Simuler la première conception cycle par cycle dans Verilator et Icarus, vérifier les deux face à un modèle Python sous une
contre-pression aléatoire, et modéliser le coût du pipeline. \textbf{État final~:} \cref{tab:nw:programmable-hardware-i-architecture-and-toolchain:cycles},
\cref{fig:nw:programmable-hardware-i-architecture-and-toolchain:pipeline} et les tests au vert de \texttt{firm.hdlkit}.
\begin{tutsteps}
\item \textbf{Lire la conception.} \texttt{hdl/hdk\_top.sv} enchaîne un tampon de débordement, l'extracteur de champ
  (\cref{lst:nw:programmable-hardware-i-architecture-and-toolchain:field}), une unité CRC-32 et le comparateur
  (\cref{lst:nw:programmable-hardware-i-architecture-and-toolchain:cmp}).
\item \textbf{Construire les deux.} \texttt{firm\_hdlkit.build\_verilator} lance \texttt{verilator --cc --exe --build} avec le \omterm{def:nw:programmable-hardware-i-architecture-and-toolchain:tb}{banc de test} en C++~;
  \texttt{build\_icarus} lance \texttt{iverilog -g2012} avec celui en SystemVerilog.
\item \textbf{Exécuter et comparer.} \texttt{write\_stim} transforme les trames en beats de 8 octets, \texttt{write\_ready} fournit la
  contre-pression de l'aval~; les deux simulateurs impriment les mêmes lignes d'événements, et \texttt{golden} calcule les champs, les CRC et les comparaisons en Python.
\item \textbf{Cycles et horloges.} \texttt{nw\_fpga.first\_frame} fournit le tableau~; \texttt{timing} et \texttt{depth\_for} évaluent
  la \cref{prop:nw:programmable-hardware-i-architecture-and-toolchain:pipe}, et \texttt{python fig\_fpga.py} écrit les données du graphique.
\end{tutsteps}
\textbf{Pour aller plus loin.} Déplacer le champ aux octets 60 à 63 et relever le cycle où il apparaît~; fixer le signal ready de l'aval à 50~\% et compter
les cycles que met l'ensemble du flux.
\end{tutorial}

\section{Construction~: la boîte à outils matérielle}\label{bld:nw:programmable-hardware-i-architecture-and-toolchain:hdlkit}

\begin{build}
\bfield{Rôle} Les primitives et le banc de vérification que chaque conception matérielle de ce livre utilise.
\bfield{Interface} modules SystemVerilog \texttt{hdk\_skid} (tampon skid, paramètre \texttt{W}), \texttt{hdk\_field} (\texttt{OFF},
\texttt{LEN}), \texttt{hdk\_crc32}, \texttt{hdk\_cmp} (\texttt{PIPE}) et \texttt{hdk\_top}~; \omterm{def:nw:programmable-hardware-i-architecture-and-toolchain:tb}{bancs de test} \texttt{tb/hdk\_tb.cpp} (Verilator)
et \texttt{tb/hdk\_tb.sv} (Icarus) qui impriment les mêmes lignes d'événements. Python \texttt{firm\_hdlkit}~: \texttt{require}, \texttt{beats},
\texttt{write\_stim}, \texttt{write\_ready}, \texttt{build\_verilator}, \texttt{build\_icarus}, \texttt{run\_verilator},
\texttt{run\_icarus}, \texttt{events}, \texttt{golden}, \texttt{CycleModel}.
\bfield{Règles} Beats de 8 octets, l'octet~0 dans la voie de poids fort~; chaque sortie enregistrée~; le HDL est accepté par Verilator 4.038 avec
\texttt{-Wall} et par Icarus 11 avec \texttt{-g2012}, sans avertissement~; une simulateur absent lève \texttt{ToolMissing}.
\bfield{Tests d'acceptation} \texttt{code/firm/hdlkit/tests/}~: les deux simulateurs identiques cycle par cycle et égaux au modèle de référence sur des
trames aléatoires avec contre-pression aléatoire, pour des comparateurs à un et deux étages~; chaque beat accepté une seule fois~; les latences en
cycles~; les fonctions beat et golden vérifiées à la main.
\bfield{Pour aller plus loin} Un FIFO asynchrone pour un \omterm{def:nw:programmable-hardware-i-architecture-and-toolchain:cdc}{franchissement de domaine d'horloge} et un test qui l'échantillonne à des horloges sans lien~; un beat plus large (512 bits).
\end{build}

\begin{omsources}
\item AMD, \emph{UltraScale Architecture Configurable Logic Block User Guide} (UG574) et \emph{Memory Resources} (UG573)~; AMD, \emph{Vitis
High-Level Synthesis User Guide} (UG1399).
\item Verilator documentation~; Icarus Verilog documentation~; IEEE~1800 (SystemVerilog).
\item J.~W. Lockwood et coauteurs, ``A Low-Latency Library in FPGA Hardware for High-Frequency Trading (HFT)'', \emph{IEEE Symposium on
High-Performance Interconnects} (2012).
\end{omsources}

\section{Exercices}

\begin{exercise}[$\star$]\label{exo:nw:programmable-hardware-i-architecture-and-toolchain:1}
Quelle horloge une conception doit-elle atteindre pour traiter une liaison de \qty{25}{\giga\bit\per\second} 32 bits à la fois, puis 128 bits à la fois~?
\end{exercise}

\begin{exercise}[$\star$]\label{exo:nw:programmable-hardware-i-architecture-and-toolchain:2}
Avec les délais du modèle, quelle est l'horloge maximale d'une décision de 12 niveaux en un seul étage, et sa latence~?
\end{exercise}

\begin{exercise}[$\star$]\label{exo:nw:programmable-hardware-i-architecture-and-toolchain:3}
Dans la \cref{tab:nw:programmable-hardware-i-architecture-and-toolchain:cycles}, à quel cycle le champ apparaîtrait-il s'il occupait les octets 12 à
15~? Et les octets 14 à 17 d'une trame de 24 octets~?
\end{exercise}

\begin{exercise}[$\star\star$]\label{exo:nw:programmable-hardware-i-architecture-and-toolchain:4}
Pourquoi la boucle du \cref{lst:nw:programmable-hardware-i-architecture-and-toolchain:field} prend-elle un cycle quelles que soient les données, alors que la
même boucle en C++ peut ne pas le faire~?
\end{exercise}

\begin{exercise}[$\star\star$]\label{exo:nw:programmable-hardware-i-architecture-and-toolchain:5}
Dans la \cref{fig:nw:programmable-hardware-i-architecture-and-toolchain:pipeline}, pourquoi six étages sont-ils plus rapides que cinq~? Qu'est-ce que cela dit
du décompte des niveaux~?
\end{exercise}

\begin{exercise}[$\star\star$]\label{exo:nw:programmable-hardware-i-architecture-and-toolchain:6}
Que fait le tampon skid qu'un registre unique ne peut pas faire, et que se passerait-il pour le flux sans lui lorsque l'aval se bloque~?
\end{exercise}

\begin{exercise}[$\star\star\star$]\label{exo:nw:programmable-hardware-i-architecture-and-toolchain:7}
\emph{Programmation.} Construire \texttt{hdk\_top} avec le champ au décalage 10 et de longueur 8 dans les deux simulateurs, exécuter les trames aléatoires du test et
vérifier les deux contre le modèle de référence. Combien de cycles après le beat qui complète la trame le champ apparaît-il~?
\end{exercise}

\begin{exercise}[$\star\star\star$]\label{exo:nw:programmable-hardware-i-architecture-and-toolchain:8}
\emph{Trouver la faille.} «~La conception passe dans Verilator, elle fonctionnera donc sur la carte~; lancer un second simulateur est redondant.~»
\end{exercise}

\section{Problème~: huit cycles}

\begin{problem}[{Problème du week-end --- une décision face à l'horloge d'une liaison}]\label{pb:nw:programmable-hardware-i-architecture-and-toolchain:1}
La décision matérielle d'une société de trading (extraire un champ, le comparer à un seuil, vérifier une limite) demande 12 niveaux de logique. Utiliser le modèle du chapitre~:
\qty{0.5}{\nano\second} par niveau, \qty{0.4}{\nano\second} par registre. La conception traite la liaison à 64 bits par cycle.

\medskip\noindent\textbf{Partie I --- Les horloges.}
\begin{enumerate}
\item Quelle horloge une liaison de \qty{10}{\giga\bit\per\second} exige-t-elle à 64 bits par cycle~? Et une liaison de \qty{25}{\giga\bit\per\second}~?
\item Quelle est l'horloge maximale de la décision en un seul étage~?
\item Respecte-t-elle l'horloge de l'une ou l'autre liaison~?
\item Combien de temps une trame de 128 octets met-elle à arriver à l'horloge de chaque liaison, en cycles puis en nanosecondes~?
\end{enumerate}

\medskip\noindent\textbf{Partie II --- Pipeline.}
\begin{enumerate}[resume]
\item Combien d'étages la liaison de \qty{25}{\giga\bit\per\second} exige-t-elle~?
\item Quelle est l'horloge maximale avec ce nombre d'étages~?
\item Quelle est la latence de la décision lorsque la conception fonctionne à l'horloge de la liaison~?
\item Quelle est la latence à l'horloge maximale propre à la conception~?
\end{enumerate}

\medskip\noindent\textbf{Partie III --- Où passe le temps.}
\begin{enumerate}[resume]
\item Si le champ se termine à l'octet 40 d'une trame de 128 octets, combien de beats doivent être arrivés pour que le champ soit complet~?
\item Ajouter la décision à l'horloge de \qty{25}{\giga\bit\per\second}~: à quel moment la décision est-elle prise, compté depuis le premier bit de la trame~?
\item Quelle part de la trame arrive encore à cet instant~?
\item Pourquoi a-t-on intérêt à décider avant l'arrivée de la séquence de contrôle de la trame, et quel en est le risque~?
\end{enumerate}

\medskip\noindent\textbf{Partie IV --- Le verdict.}
\begin{enumerate}[resume]
\item Énoncer le \emph{résultat nommé}~: la profondeur de pipeline qui respecte l'horloge de \qty{25}{\giga\bit\per\second} et la latence de la décision, face
  à la conception sans pipeline à sa propre horloge.
\item Un chemin de données de 128 bits changerait-il la réponse~?
\item Que remplacerait un rapport de temporisation réel dans ce modèle~?
\item Pourquoi le pipeline ne réduit-il jamais la latence~?
\item Pourquoi le teneur de marché accepterait-il une latence de décision plus longue pour une liaison plus rapide~?
\item Comment vérifieriez-vous la conception pipelinée face à la conception sans pipeline~?
\item Qu'ajoute un \omterm{def:nw:programmable-hardware-i-architecture-and-toolchain:cdc}{franchissement de domaine d'horloge} si la décision s'exécute à une horloge autre que celle de la liaison~?
\item En une phrase~: que décide le débit de la liaison sur la conception~?
\end{enumerate}
\end{problem}

\section{Questions d'entretien}

\begin{interviewq}[$\star$ \iqroles{développeur}]\label{iq:nw:programmable-hardware-i-architecture-and-toolchain:1}
Qu'est-ce qu'un FPGA, et pourquoi les sociétés de trading en utilisent-elles~?
\end{interviewq}

\begin{interviewq}[$\star\star$ \iqroles{développeur}]\label{iq:nw:programmable-hardware-i-architecture-and-toolchain:2}
Qu'est-ce que la \omterm{def:nw:programmable-hardware-i-architecture-and-toolchain:timing}{fermeture temporelle}~(timing closure)~? Que faites-vous lorsqu'un chemin échoue à la temporisation~?
\end{interviewq}

\begin{interviewq}[$\star\star$ \iqroles{développeur}]\label{iq:nw:programmable-hardware-i-architecture-and-toolchain:3}
Le pipeline ajoute des cycles. Comment peut-il rendre une conception plus rapide~?
\end{interviewq}

\begin{interviewq}[$\star\star$ \iqroles{développeur}]\label{iq:nw:programmable-hardware-i-architecture-and-toolchain:4}
RTL ou \omterm{def:nw:programmable-hardware-i-architecture-and-toolchain:hls}{synthèse de haut niveau} pour une conception de trading~? Argumentez les deux positions.
\end{interviewq}

\begin{interviewq}[$\star\star\star$ \iqroles{développeur}]\label{iq:nw:programmable-hardware-i-architecture-and-toolchain:5}
Comment vérifieriez-vous un analyseur de flux matériel avant qu'il ne fonctionne face à une bourse en direct~?
\end{interviewq}

\begin{interviewq}[$\star\star$ \iqroles{développeur}]\label{iq:nw:programmable-hardware-i-architecture-and-toolchain:6}
Qu'est-ce qu'un \omterm{def:nw:programmable-hardware-i-architecture-and-toolchain:cdc}{franchissement de domaine d'horloge}, et qu'est-ce qui va de travers s'il n'est pas géré~?
\end{interviewq}
